\subsection*{把复杂周期形状拆成简单运动}
一条复杂周期曲线，可以尝试写成许多不同频率的正弦、余弦之和。每个频率只承担一种规则振动，幅度与相位说明它贡献多少、从哪里开始。复指数把一对正弦余弦合成平面旋转，使正负频率能用统一记号表达。

关键问题是怎样测出某个频率的系数。把信号乘上反向旋转的目标频率：与目标同频的部分变成常量，其他整倍频仍在一周期内转完若干圈，平均为零。因此积分能够挑出匹配部分，这就是正交性的作用。

\subsection*{系数公式之后还有收敛问题}
对 $3+2\cos t$，常量三不旋转，余弦可拆成两个幅度为一、方向相反的旋转。这是可逐项验证的有限例子。先把它算明白，再处理一般系数公式，会更容易看清归一化因子和正负频率。

在热实验中改变高频振幅，可以看到同一个低频骨架怎样叠加出细节；随后推进时间，高频先衰减。傅里叶分解因而不仅是换一种画法，还把原本耦合的演化拆成独立系数的变化。

并非写出形式级数就证明它处处等于原函数。对跳跃点，常见分段光滑条件下级数趋向左右极限平均，还会出现局部振铃。讨论近似时必须说明点态、均方或其他意义。这里先掌握频率测量机制，把完备性和收敛留作明确的分析问题。
